\documentclass[11pt,a4paper]{article}

\usepackage[utf8]{inputenc}
\usepackage[spanish,es-nodecimaldot,es-noquoting]{babel}
\usepackage[margin=2cm]{geometry}
\usepackage{booktabs}
\usepackage{tabularx}
\usepackage{array}
\usepackage{xcolor}
\usepackage{hyperref}
\usepackage{xurl}
\usepackage{amsmath,amssymb}
\usepackage{enumitem}
\usepackage{tcolorbox}
\usepackage{titlesec}
\usepackage{fancyhdr}
\usepackage{colortbl}
\usepackage{longtable}
\usepackage{tikz}
\usetikzlibrary{shapes.geometric, arrows.meta, positioning}

% Paleta de colores formal
\definecolor{primary}{RGB}{30, 64, 175}
\definecolor{secondary}{RGB}{71, 85, 105}
\definecolor{accentgreen}{RGB}{22, 101, 52}
\definecolor{accentred}{RGB}{153, 27, 27}
\definecolor{lightgray}{RGB}{241, 245, 249}
\definecolor{tableheader}{RGB}{226, 232, 240}

\hypersetup{
    colorlinks=true,
    linkcolor=primary,
    filecolor=primary,
    urlcolor=primary,
    citecolor=primary
}

\emergencystretch 3em
\hfuzz 2pt

\pagestyle{fancy}
\fancyhf{}
\fancyhead[L]{\small \textcolor{secondary}{Version-Eco: Plan Maestro de Pruebas y Reporte de Ejecución}}
\fancyhead[R]{\small \textcolor{secondary}{Javeriana Cali -- 2026}}
\fancyfoot[C]{\thepage}
\renewcommand{\headrulewidth}{0.4pt}

\titleformat{\section}{\Large\bfseries\color{primary}}{\thesection.}{0.5em}{}[\titlerule]
\titleformat{\subsection}{\large\bfseries\color{secondary}}{\thesubsection.}{0.5em}{}
\titleformat{\subsubsection}{\normalsize\bfseries\color{secondary}}{\thesubsubsection.}{0.5em}{}

\begin{document}

% ==============================================================================
% PORTADA
% ==============================================================================
\begin{titlepage}
    \centering
    \vspace*{0.5cm}
    
    {\Huge \textbf{Pontificia Universidad Javeriana Cali}}\\[0.3cm]
    {\LARGE Facultad de Ingeniería y Ciencias}\\[0.2cm]
    {\Large Departamento de Ingeniería de Sistemas y Computación}\\[1.5cm]
    
    {\huge \bfseries \color{primary} Plan Maestro de Pruebas, Diseño de Casos (Caja Negra) y Reporte de Ejecución}\\[0.5cm]
    {\LARGE \textbf{Sistema de Gestión de Inventario y Producción}}\\[0.2cm]
    {\Large \textbf{Empresa Version-Eco}}\\[1.8cm]
    
    \textbf{Asignatura:}\\
    Construcción de Software y Pruebas\\[1.2cm]
    
    \textbf{Integrantes del Equipo:}\\
    Harold Guerrero\\
    Benjamin Ortiz\\
    Samuel Rueda\\
    Daniela Orobio\\[1.5cm]
    
    \vfill
    {\large Santiago de Cali, Colombia\\ 21 de octubre de 2026}
\end{titlepage}

\tableofcontents
\newpage

% ==============================================================================
% 1. INFORMACIÓN GENERAL DEL PROYECTO
% ==============================================================================
\section{Información General del Proyecto}

\begin{itemize}[leftmargin=1.5cm]
    \item \textbf{Nombre del Software:} Sistema de Gestión de Inventario y Producción para Version-Eco.
    \item \textbf{Tipo de Sistema:} Aplicación Web Transaccional de Gestión de Cadena de Suministro (SMR/ERP ligero para formulación, producción centralizada y control multisede).
    \item \textbf{Integrantes del Equipo:}
    \begin{itemize}
        \item Harold Guerrero
        \item Benjamin Ortiz
        \item Samuel Rueda
        \item Daniela Orobio
    \end{itemize}
    \item \textbf{URL del Repositorio de Código:} \url{https://github.com/Benjamin-Ortiz/Version_EcoProject}
    \item \textbf{URL de la Aplicación:}
    \begin{itemize}
        \item \textbf{Ambiente Local de Pruebas (SUT):} Frontend en \url{http://localhost:5173} y Backend API en \url{http://localhost:4000/api}.
        \item \textbf{Ambiente de Despliegue en la Nube:} Despliegue en plataforma Vercel con integración a base de datos PostgreSQL Serverless (Neon).
    \end{itemize}
\end{itemize}

% ==============================================================================
% 2. DESCRIPCIÓN DEL SISTEMA
% ==============================================================================
\section{Descripción del Sistema}

\subsection{¿Qué problema resuelve el software?}
La fundación y empresa \textbf{Version-Eco} elabora productos cosméticos y de cuidado personal con enfoque vegano y ecológico (tales como champús sólidos, desodorantes naturales en crema y protectores solares minerales biodegradables). Históricamente, la empresa controlaba sus existencias mediante archivos dispersos de Microsoft Excel, lo cual generaba:
\begin{enumerate}
    \item Desconexión entre los insumos adquiridos y el gasto real de ingredientes en las órdenes de producción.
    \item Quiebres imprevistos de materias primas críticas que paralizaban la planta de elaboración.
    \item Carencia de visibilidad consolidada sobre el stock existente en cada una de sus sedes comerciales periféricas.
    \item Discrepancias en el arqueo diario al registrar ventas o mermas manualmente sin una bitácora auditada.
\end{enumerate}
El sistema resuelve esta problemática centralizando todo el ciclo operativo: adquisición de materias primas, configuración de recetas industriales (BOM), producción con descuento atómico, traslado a sucursales y generación de alertas tempranas de desabastecimiento.

\subsection{¿Quiénes son los usuarios del sistema?}
\begin{itemize}
    \item \textbf{Jefe de Producción / Maestro Formulador:} Administra las fórmulas oficiales, define las cantidades requeridas por lote y autoriza las órdenes de manufactura en la sede central.
    \item \textbf{Encargado de Bodega y Compras:} Registra el ingreso de materias primas, valida los niveles de seguridad y gestiona los despachos a puntos de venta.
    \item \textbf{Administradores de Sede:} Consultan el inventario físico disponible en su sucursal y reportan ajustes diarios de inventario derivados de ventas locales.
\end{itemize}

\subsection{¿Qué valor ofrece la aplicación?}
\begin{itemize}
    \item \textbf{Integridad y Atomicidad Operativa:} Garantiza matemáticamente que no se fabrique ningún lote si falta una sola fracción de materia prima, evitando deducciones parciales erróneas.
    \item \textbf{Trazabilidad Completa:} Cada transformación física de insumo a producto terminado genera movimientos inmutables con fecha, sede y referencia.
    \item \textbf{Toma de Decisiones Proactiva:} Semáforos de stock bajo que advierten oportunamente cuándo reabastecer insumos o trasladar productos.
\end{itemize}

% ==============================================================================
% 3. ARQUITECTURA GENERAL
% ==============================================================================
\section{Arquitectura General del Sistema}

\subsection{Tecnologías Utilizadas}
\begin{itemize}
    \item \textbf{Frontend (Capa de Cliente):} Desarrollado como una Single Page Application (SPA) con \textbf{React 18} y el empaquetador \textbf{Vite}. Maneja el enrutamiento con \textbf{React Router DOM v6}, peticiones HTTP asíncronas con \textbf{Axios} e iconografía mediante \textbf{Lucide React}.
    \item \textbf{Backend (Capa de Negocio y API):} Desarrollado en \textbf{Node.js (v20+)} utilizando \textbf{Express 4} con soporte nativo de módulos ECMAScript (ESM). Emplea \textbf{Morgan} para registro de peticiones, \textbf{CORS} y \textbf{PDFKit} para generación en tiempo real de fichas técnicas descargables en PDF.
    \item \textbf{Base de Datos (Capa de Persistencia):} Motor relacional \textbf{PostgreSQL 18}. Utiliza un pool de conexiones optimizado (\texttt{pg.Pool}), vistas calculadas (\texttt{vw\_alertas\_*}, \texttt{vw\_inventario\_central}), disparadores de auditoría y procedimientos almacenados en \textbf{PL/pgSQL} para ejecutar la lógica transaccional crítica.
\end{itemize}

\subsection{Componentes Principales del Sistema}
\begin{enumerate}
    \item \textbf{Interfaz Web SPA:} Compuesta por módulos reactivos: Dashboard de Métricas, Catálogo de Materias Primas, Catálogo de Productos y Fórmulas, Gestión de Sedes, Módulo de Producción Central, Módulo de Distribución y Vista de Movimientos.
    \item \textbf{API RESTful:} Enrutador centralizado que expone endpoints bajo el prefijo \texttt{/api/}, gestiona códigos de estado HTTP (200, 201, 400, 404, 409, 500) y normaliza excepciones de base de datos.
    \item \textbf{Motor de Reglas de Base de Datos (PL/pgSQL):} Encapsula la lógica transaccional de producción (\texttt{fn\_registrar\_produccion}) y distribución (\texttt{fn\_distribuir\_desde\_central}).
\end{enumerate}

\newpage
\subsection{Diagrama de Arquitectura}

\begin{center}
\begin{tikzpicture}[
    layerbox/.style={
        rectangle, rounded corners=6pt, line width=1.2pt,
        text width=13.5cm, align=center, inner sep=6pt
    },
    conn/.style={
        {Latex[length=3mm,width=2.2mm]}-{Latex[length=3mm,width=2.2mm]},
        line width=1.4pt
    },
    labelnode/.style={
        fill=white, draw=gray!40, rounded corners=3pt, font=\scriptsize\bfseries,
        inner sep=3pt, text=secondary
    }
]

% Nodo Capa 1: Frontend
\node[layerbox, draw=primary, fill=blue!4] (frontend) {
    {\large \textbf{\color{primary} Capa 1: Presentación y Cliente Web (Frontend SPA)}}\\[1mm]
    \textbf{React 18} \quad $\bullet$ \quad \textbf{Vite} \quad $\bullet$ \quad \textbf{React Router DOM v6} \quad $\vert$ \quad \texttt{http://localhost:5173}\\[1.5mm]
    \footnotesize
    \textbf{Módulos de Usuario:} Dashboard, Materias Primas, Productos, Fórmulas, Producción, Traslados\\
    \textbf{Servicios UI:} Cliente HTTP Axios, Manejo de Estado Local, Validación de Entradas, Lucide Icons
};

% Nodo Capa 2: Backend
\node[layerbox, draw=secondary, fill=gray!6, below=1.2cm of frontend] (backend) {
    {\large \textbf{\color{secondary} Capa 2: Lógica de Negocio y Servicios (Backend REST API)}}\\[1mm]
    \textbf{Node.js 20} \quad $\bullet$ \quad \textbf{Express 4 (ESM)} \quad $\bullet$ \quad \textbf{PDFKit} \quad $\vert$ \quad \texttt{http://localhost:4000/api}\\[1.5mm]
    \footnotesize
    \textbf{Enrutador de Endpoints:} \texttt{/api/produccion}, \texttt{/api/inventario}, \texttt{/api/sedes}, \texttt{/api/materias-primas}\\
    \textbf{Servicios:} Middlewares (CORS, Morgan, ErrorHandler), Generador PDF de Recetas, Pool TCP \texttt{pg}
};

% Nodo Capa 3: Base de Datos
\node[layerbox, draw=accentgreen, fill=green!4, below=1.2cm of backend] (database) {
    {\large \textbf{\color{accentgreen} Capa 3: Persistencia y Reglas de Negocio (Base de Datos Relacional)}}\\[1mm]
    \textbf{PostgreSQL 18} \quad $\vert$ \quad \texttt{127.0.0.1:5432 / version\_eco\_db}\\[1.5mm]
    \footnotesize
    \textbf{Tablas \& Triggers:} \texttt{materia\_prima}, \texttt{producto}, \texttt{inventario\_central}, \texttt{inventario\_sede}, \texttt{trg\_historial}\\
    \textbf{Lógica Transaccional PL/pgSQL:} \texttt{fn\_registrar\_produccion}, \texttt{fn\_distribuir\_desde\_central}, Vistas de Alerta
};

% Flechas con etiquetas descriptivas
\draw[conn, draw=primary] (frontend) -- (backend) 
    node[midway, labelnode] {PETICIONES HTTP REST / RESPUESTAS ASÍNCRONAS JSON};

\draw[conn, draw=secondary!90!black] (backend) -- (database) 
    node[midway, labelnode] {CONEXIÓN TCP (DRIVER pg) / SENTENCIAS SQL \& TRANSACCIONES ACID};

\end{tikzpicture}\\[2mm]
{\small \textbf{Figura 1:} Diagrama de Arquitectura en Tres Capas (C4 - Nivel de Contenedores) para Version-Eco.}
\end{center}

% ==============================================================================
% 4. ACTORES O TIPOS DE USUARIO
% ==============================================================================
\section{Actores o Tipos de Usuario}

\begin{table}[htbp]
\centering
\small
\begin{tabularx}{\textwidth}{p{3.5cm} p{3.5cm} X}
\toprule
\textbf{Actor} & \textbf{Nivel de Acceso} & \textbf{Responsabilidad en el Sistema} \\
\midrule
\textbf{Jefe de Producción} & Operación Completa (\textit{Master}) & Crea y edita recetas/fórmulas, programa órdenes de fabricación, valida existencias globales y autoriza traslados a sucursales. \\
\addlinespace
\textbf{Encargado de Almacén} & Operativo / Bodega & Registra el ingreso de nuevos insumos, abastece existencias físicas y atiende las alertas rojas de materias primas agotadas. \\
\addlinespace
\textbf{Operario de Sede} & Consulta y Cuadre Local & Monitorea el stock disponible en su sede asignada y aplica ajustes manuales de inventario por ventas diarias o mermas. \\
\bottomrule
\end{tabularx}
\caption{Matriz de Actores y Roles del Sistema.}
\end{table}

\textit{Nota: En la versión 1.0 del aplicativo, el sistema opera con interfaz abierta basada en rol único operativo (Master), facilitando el acceso rápido en estaciones de planta.}

% ==============================================================================
% 5. LISTA DE FUNCIONALIDADES DEL SISTEMA (NO CRUD)
% ==============================================================================
\section{Lista de Funcionalidades Complejas del Sistema (No CRUD)}

Siguiendo la selección formal acordada por el equipo, las cuatro funcionalidades identificadas que contienen lógica de negocio compleja son:

\subsection{Funcionalidad F1: Gestión de Materia Prima}
\begin{itemize}
    \item \textbf{Actor que la utiliza:} Encargado de Almacén y Compras.
    \item \textbf{Descripción:} Administra el ciclo de vida de los ingredientes con reglas de validación compuestas:
    \begin{itemize}
        \item Control de unicidad de nombre sensible a espacios y mayúsculas (RN-01).
        \item Control de rango estricto: la cantidad debe ser $\ge 0$ (RN-02) y no puede superar el tope máximo de $50$ unidades (RN-04).
        \item Restricción de unidad de medida predefinida (\texttt{kg}, \texttt{L}, \texttt{g}, etc.) (RN-05).
        \item Bloqueo de eliminación física si la materia prima pertenece a una fórmula de producción activa (RN-06).
        \item Historial automático de auditoría ante cualquier cambio en cantidad o datos básicos (RN-07).
    \end{itemize}
    \item \textbf{Ejemplo:} Si se intenta registrar \textit{"Aceite de coco"} con cantidad $50.01\,\text{kg}$, el sistema lo rechaza por superar el límite de almacenamiento de 50. Si se intenta eliminar \textit{"Gel de aloe vera"} estando vinculada a la fórmula del champú, el sistema bloquea la acción.
\end{itemize}

\subsection{Funcionalidad F2: Creación y Definición de Fórmulas}
\begin{itemize}
    \item \textbf{Actor que la utiliza:} Jefe de Producción.
    \item \textbf{Descripción:} Permite estructurar la lista de materiales (Bill of Materials) requerida para la fabricación de cada producto:
    \begin{itemize}
        \item Valida que el producto tenga un nombre único en el catálogo.
        \item Exige que la fórmula contenga al menos una materia prima existente (RN-08).
        \item Exige que la cantidad requerida de cada insumo sea estrictamente mayor que cero ($>0$).
        \item Bloquea el ingreso redundante de una misma materia prima dentro de la misma receta (RN-14).
    \end{itemize}
    \item \textbf{Ejemplo:} Al crear la fórmula del Champú, si se intenta asociar dos veces el insumo \textit{"Aloe Vera"} o se ingresa cantidad $0\,\text{kg}$, el sistema rechaza la definición de la fórmula.
\end{itemize}

\subsection{Funcionalidad F3: Producción en Sede Central con Consumo Atómico}
\begin{itemize}
    \item \textbf{Actor que la utiliza:} Jefe de Producción.
    \item \textbf{Descripción:} Registra la elaboración de $Q$ unidades de producto a partir de su fórmula, ejecutando un consumo atómico de materias primas:
    \begin{itemize}
        \item Calcula el consumo necesario: $Consumo_i = RequeridoUnitario_i \times Q$.
        \item Si falta stock de \textbf{cualquier} insumo, se bloquea la orden completa (\textit{Rollback}), sin descontar ningún otro ingrediente (RN-10).
        \item Si todas las materias primas alcanzan, se descuentan simultáneamente, se incrementa el inventario central del producto terminado y se genera asiento inmutable en la bitácora de movimientos (RN-09).
        \item Exige que la cantidad $Q$ sea un entero positivo mayor que cero (RN-12, RN-15).
    \end{itemize}
    \item \textbf{Ejemplo:} Para fabricar 10 champús que requieren $5\,\text{kg}$ de Aloe y $2\,\text{L}$ de Aceite, si en bodega hay $10\,\text{kg}$ de Aloe pero solo $1.95\,\text{L}$ de Aceite, el sistema aborta la operación y no descuenta nada del Aloe.
\end{itemize}

\subsection{Funcionalidad F4: Traslado de Productos a Sedes}
\begin{itemize}
    \item \textbf{Actor que la utiliza:} Jefe de Producción / Encargado de Logística.
    \item \textbf{Descripción:} Gestiona el envío físico de productos terminados desde la bodega central hacia las sedes comerciales periféricas:
    \begin{itemize}
        \item Valida que el origen sea exclusivamente la sede central y que la cantidad $Q$ sea un entero mayor que cero (RN-11, RN-12).
        \item Comprueba que la cantidad a despachar no exceda las existencias disponibles en el inventario central.
        \item Deduce el saldo de central, suma exactamente la misma cantidad al inventario de la sede destino y genera dos registros cruzados de auditoría (\texttt{DISTRIBUCION\_SALIDA} y \texttt{DISTRIBUCION\_ENTRADA}).
    \end{itemize}
    \item \textbf{Ejemplo:} Si la central tiene 20 unidades y se solicita enviar 21 unidades a la Sede Norte, el sistema rechaza la solicitud. Si se envían 5 unidades, la central queda con 15 y la Sede Norte con 5.
\end{itemize}

% ==============================================================================
% 6. REGLAS DE NEGOCIO IDENTIFICADAS
% ==============================================================================
\section{Reglas de Negocio Identificadas}

\begin{longtable}{p{1.8cm} p{10.8cm} p{3.2cm}}
\caption{Matriz Completa de Reglas de Negocio del Sistema Version-Eco.}\label{tab:reglas}\\
\toprule
\textbf{ID Regla} & \textbf{Enunciado de la Regla de Negocio} & \textbf{Origen / Requisito} \\
\midrule
\endfirsthead
\multicolumn{3}{c}{\tablename\ \thetable\ -- \textit{Continuación de la página anterior}} \\
\toprule
\textbf{ID Regla} & \textbf{Enunciado de la Regla de Negocio} & \textbf{Origen / Requisito} \\
\midrule
\endhead
\midrule
\multicolumn{3}{r}{\textit{Continúa en la siguiente página\dots}} \\
\endfoot
\bottomrule
\endlastfoot
\textbf{RN-01} & El nombre de materia prima y producto es único. La comparación ignora mayúsculas y espacios sobrantes (``Aloe'' equivale a ``aloe ''). & RF-01, RF-04, RNF-03 \\
\addlinespace
\textbf{RN-02} & La cantidad de materia prima puede ser 0, pero nunca negativa. Si queda en 0 se genera una alerta. & Aclaración del equipo \\
\addlinespace
\textbf{RN-03} & Ninguna acción permitida puede dejar una materia prima o un inventario con cantidad negativa. & Aclaración del equipo \\
\addlinespace
\textbf{RN-04} & La cantidad de una materia prima no puede superar 50 unidades de su medida ($50\,\text{kg}, 50\,\text{L}$). Cualquier alta o ingreso que la deje por encima de 50 se rechaza. & Aclaración del equipo \\
\addlinespace
\textbf{RN-05} & La unidad de medida se selecciona de una lista predefinida (\texttt{kg}, \texttt{L}, \texttt{g}, \texttt{ml}, \texttt{unidad}). & RF-01 \\
\addlinespace
\textbf{RN-06} & Una materia prima asociada a una fórmula no se puede eliminar. La eliminación exige confirmación. & RF-03 \\
\addlinespace
\textbf{RN-07} & Toda modificación de una materia prima genera un registro de historial. & RF-02 \\
\addlinespace
\textbf{RN-08} & Una fórmula debe tener al menos una materia prima existente con cantidad $>0$. No puede repetirse dentro de la misma fórmula. & RF-05 \\
\addlinespace
\textbf{RN-09} & Consumo de materia prima en una producción = cantidad unitaria de fórmula $\times$ cantidad producida. & RF-09 \\
\addlinespace
\textbf{RN-10} & Si alguna materia prima es insuficiente, no se registra la producción y no se descuenta nada (operación atómica). & RF-09 \\
\addlinespace
\textbf{RN-11} & Solo se puede trasladar desde la sede central y la cantidad no puede superar el stock disponible en ella. & Aclaración del equipo \\
\addlinespace
\textbf{RN-12} & La cantidad a producir o trasladar debe ser un número entero mayor que 0. & Aclaración del equipo \\
\addlinespace
\textbf{RN-13} & Solo existe el rol \textit{master}; no hay inicio de sesión en esta versión. & Aclaración del equipo \\
\addlinespace
\textbf{RN-14} & Si se intenta agregar la misma materia prima dos veces a una fórmula, el sistema debe rechazarlo. & Deriva de RN-08 \\
\addlinespace
\textbf{RN-15} & Producir o trasladar con cantidades decimales, 0 o negativas se rechaza. & Deriva de RN-12 \\
\addlinespace
\textbf{RN-16} & El inventario de un producto puede llegar hasta 100 unidades (esta regla no limita la producción por sí misma). & Aclaración del equipo \\
\end{longtable}

% ==============================================================================
% 7. DOCUMENTO DE PLAN DE PRUEBAS (ESTRUCTURA COMPLETA)
% ==============================================================================
\section{Documento de Plan Maestro de Pruebas (PMP)}

\subsection{Introducción}
El presente plan de pruebas establece las directrices para verificar el Sistema de Gestión de Inventario para Version-Eco. El sistema tiene como finalidad mejorar el control de materias primas, productos terminados, fórmulas de producción, sedes e inventarios. Además, permite registrar procesos de producción, descontar automáticamente las materias primas utilizadas, realizar ajustes manuales de inventario y generar alertas cuando los niveles de existencia sean bajos.

Las pruebas están orientadas a verificar que las funcionalidades cumplan con los requisitos definidos y que el sistema mantenga la integridad y consistencia de la información durante las operaciones de los usuarios.

\subsection{Alcance de las Pruebas}
\subsubsection{Funcionalidades que serán probadas}
El plan contempla las funcionalidades principales definidas en el levantamiento de requisitos:
\begin{itemize}
    \item \textbf{RF-01:} Registro de materias primas.
    \item \textbf{RF-02:} Modificación de materias primas.
    \item \textbf{RF-03:} Eliminación de materias primas.
    \item \textbf{RF-04:} Registro de productos.
    \item \textbf{RF-05:} Definición de fórmulas de producción.
    \item \textbf{RF-06:} Registro y gestión de sedes.
    \item \textbf{RF-07:} Consulta de inventario por sede.
    \item \textbf{RF-08:} Registro de producción.
    \item \textbf{RF-09:} Descuento automático de materias primas.
    \item \textbf{RF-10:} Ajuste manual de inventario.
    \item \textbf{RF-11:} Alerta de stock mínimo de materias primas.
    \item \textbf{RF-12:} Alerta de stock mínimo de productos.
    \item Requisitos no funcionales: usabilidad (RNF-01), disponibilidad (RNF-02), integridad (RNF-03) y rendimiento (RNF-04).
\end{itemize}

\subsubsection{Fuera del alcance}
\begin{itemize}
    \item Pruebas de sistemas externos utilizados para registrar ventas comerciales.
    \item Pruebas de infraestructura física o de hardware de la empresa.
    \item Pruebas de facturación, pasarelas de pago o contabilidad financiera.
    \item Pruebas de aplicaciones móviles (el alcance es estrictamente web).
\end{itemize}

\subsection{Criterios de Entrada y Salida}
\subsubsection{Criterios de Entrada -- Definition of Ready (DoR)}
Una funcionalidad podrá comenzar su etapa de pruebas cuando cumpla las siguientes 10 condiciones:
\begin{enumerate}
    \item El requisito funcional se encuentra identificado y documentado.
    \item La funcionalidad ha sido desarrollada.
    \item La aplicación puede ejecutarse sin errores críticos.
    \item La base de datos se encuentra disponible.
    \item El backend puede conectarse correctamente con PostgreSQL.
    \item El frontend puede comunicarse con la API.
    \item Existe información o datos de prueba suficientes para ejecutar los casos.
    \item La funcionalidad se encuentra disponible en el ambiente de pruebas.
    \item Se conocen los resultados esperados de cada operación.
    \item Las dependencias necesarias para la funcionalidad se encuentran implementadas (por ejemplo, para probar fórmulas deben existir materias primas y productos registrados).
\end{enumerate}

\subsubsection{Criterios de Salida -- Definition of Done (DoD)}
Una funcionalidad se considerará aprobada cuando:
\begin{enumerate}
    \item Todos los casos de prueba críticos hayan sido ejecutados.
    \item Los casos relacionados con requisitos de prioridad alta hayan sido aprobados.
    \item No existan defectos críticos abiertos.
    \item Los errores encontrados hayan sido documentados formalmente.
    \item Los defectos corregidos hayan sido probados nuevamente.
    \item Se hayan realizado pruebas de regresión sobre las funcionalidades relacionadas.
    \item Los datos almacenados en la base de datos correspondan con las operaciones realizadas.
    \item Se hayan registrado las evidencias de ejecución.
    \item Se haya elaborado el informe de resultados.
    \item El equipo considere que la funcionalidad cumple con los requisitos definidos.
\end{enumerate}

\subsubsection{Criterio de Aceptación Global}
El proyecto se considera apto para entrega cuando el \textbf{100\% de los casos de prioridad alta} hayan sido ejecutados y aprobados, al menos el \textbf{90\% del total de casos} resulten exitosos, y no existan defectos críticos ni de severidad alta abiertos al momento del cierre.

\subsection{Estrategia de Pruebas}
Se implementan 10 tipos de pruebas combinadas:
\begin{itemize}
    \item \textbf{Pruebas Funcionales Positivas:} Comprobar que el sistema permita completar operaciones con datos válidos.
    \item \textbf{Pruebas Funcionales Negativas:} Provocar intentos no permitidos (duplicados, cantidades inválidas, insumos insuficientes).
    \item \textbf{Pruebas de Integración:} Evaluar procesos continuos (producción $\rightarrow$ descuento $\rightarrow$ inventario central $\rightarrow$ traslado).
    \item \textbf{Pruebas de API:} Comprobar endpoints REST y respuestas de estado HTTP.
    \item \textbf{Pruebas de Interfaz (GUI):} Validar formularios, campos obligatorios, botones y alertas visuales.
    \item \textbf{Pruebas de Base de Datos:} Verificar que los registros, llaves foráneas y movimientos coincidan en PostgreSQL.
    \item \textbf{Pruebas de Regresión:} Volver a ejecutar casos tras correcciones en producción o inventario.
    \item \textbf{Pruebas de Rendimiento:} Validar que las consultas de inventario tomen menos de 2 segundos (RNF-04).
    \item \textbf{Pruebas de Usabilidad:} Verificar que el registro de materias primas tome menos de 2 minutos (RNF-01).
\end{itemize}

\subsection{Entorno y Configuración de Pruebas}
\begin{itemize}
    \item \textbf{Sistema Operativo:} Windows 10/11.
    \item \textbf{Base de Datos:} PostgreSQL 18 ejecutándose localmente en puerto 5432 (\texttt{version\_eco\_db}).
    \item \textbf{Backend:} Node.js 20+, Express 4 en \texttt{http://localhost:4000/api}.
    \item \textbf{Frontend:} React 18, Vite en \texttt{http://localhost:5173}.
    \item \textbf{Navegador:} Google Chrome / Mozilla Firefox actualizado.
\end{itemize}

\subsection{Gestión de Riesgos y Clasificación de Defectos}
\begin{table}[htbp]
\centering
\small
\begin{tabularx}{\textwidth}{p{4.2cm} c c X}
\toprule
\textbf{Riesgo} & \textbf{Probabilidad} & \textbf{Impacto} & \textbf{Acción de Mitigación} \\
\midrule
Fallo de conexión con PostgreSQL & Media & Alto & Verificar configuración y script antes de las pruebas. \\
Datos de prueba insuficientes & Media & Medio & Cargar scripts de seed demo previamente. \\
Defectos en producción que afectan inventario & Media & Alto & Priorizar pruebas de producción y descuentos. \\
Cálculos incorrectos de recetas & Media & Alto & Comparar fórmulas con cantidades resultantes en BD. \\
Inventarios negativos & Baja & Alto & Validar restricciones \texttt{CHECK} a nivel relacional. \\
Eliminación accidental de información & Baja & Alto & Utilizar base de datos de pruebas aislada. \\
Diferencias entre frontend y backend & Media & Alto & Ejecutar pruebas de integración completas. \\
\bottomrule
\end{tabularx}
\caption{Matriz de Riesgos del Plan de Pruebas.}
\end{table}

\textbf{Clasificación de Defectos:}
\begin{itemize}
    \item \textbf{Crítico:} Impide utilizar una función esencial o provoca pérdida/corrupción de información.
    \item \textbf{Alto:} Afecta considerablemente una funcionalidad principal sin alternativa inmediata.
    \item \textbf{Medio:} Existe el error, pero se puede continuar operando el sistema.
    \item \textbf{Bajo:} Problemas cosméticos o menores que no impiden el funcionamiento.
\end{itemize}

\subsection{Roles y Responsabilidades}
\begin{itemize}
    \item \textbf{Líder de pruebas (Benjamin Ortiz):} Coordinar la ejecución y controlar el cumplimiento del plan.
    \item \textbf{Analista de pruebas (Samuel Rueda):} Diseñar casos y relacionarlos con los requisitos y técnicas de caja negra.
    \item \textbf{Ejecutor de pruebas (Harold Guerrero):} Ejecutar casos en el navegador y almacenar evidencias.
    \item \textbf{Responsable de defectos (Daniela Orobio):} Registrar los errores, documentar causa raíz y realizar seguimiento.
\end{itemize}

\subsection{Observación de Trazabilidad entre Requisito e Implementación}
En la documentación inicial, el requisito RF-08 indicaba que al registrar producción se especificaba la sede donde se elaboró. Sin embargo, la arquitectura actual maneja primero un \textbf{Inventario Central} donde ingresa el producto terminado y, posteriormente, una operación independiente de \textbf{Distribución} hacia las sedes. Esta diferencia se documenta formalmente para mantener la trazabilidad entre el requisito de negocio y la implementación real.

% ==============================================================================
% 8. DISEÑO DE CASOS DE PRUEBA (CAJA NEGRA)
% ==============================================================================
\section{Diseño de Casos de Prueba (Caja Negra)}

\subsection{Relación entre Funcionalidades y Casos de Prueba (Cobertura Intrínseca)}
En el contexto de aseguramiento de calidad, existe una relación jerárquica y de dependencia directa entre las \textbf{Funcionalidades (Features)} y los \textbf{Casos de Prueba (Test Cases)}:
\begin{itemize}
    \item \textbf{Las Funcionalidades (F1, F2, F3, F4)} representan los macro-procesos del negocio que agrupan y satisfacen los requisitos funcionales (RF) del cliente.
    \item \textbf{Los Casos de Prueba (31 casos)} son las ejecuciones atómicas diseñadas mediante técnicas formales (Tablas de Decisión, Arreglos Ortogonales y Valores Límite) que someten a prueba las entradas, reglas de validación y límites de cada funcionalidad.
    \item \textbf{Interdependencia y Cobertura Intrínseca (Cadena de Valor):}\\
    El sistema no opera con módulos aislados, sino con una cadena de custodia del inventario:
    \begin{equation*}
        \text{F1 (Materias Primas)} \xrightarrow{\text{Insumo para}} \text{F2 (Fórmulas)} \xrightarrow{\text{Receta para}} \text{F3 (Producción Central)} \xrightarrow{\text{Stock para}} \text{F4 (Traslado a Sedes)}
    \end{equation*}
    Por tanto, \textbf{al probar una funcionalidad aguas abajo se evalúan intrínsecamente las funcionalidades previas}:
    \begin{enumerate}[noitemsep]
        \item Para definir una fórmula (\textbf{F2}), es un prerrequisito obligatorio haber creado los insumos en \textbf{F1}.
        \item Al ejecutar una orden de manufactura (\textbf{F3}), el sistema valida intrínsecamente la fórmula (\textbf{F2}), descuenta el stock de insumos (\textbf{F1}), actualiza el inventario central y dispara reactivamente las alertas en el Dashboard.
        \item Para realizar una distribución a sucursales (\textbf{F4}), se requiere forzosamente el producto terminado fabricado en \textbf{F3} y las sedes registradas en el sistema.
    \end{enumerate}
\end{itemize}

\subsection{Prerrequisitos de la Base de Datos y Datos de Prueba (Baseline)}
Para garantizar la repetibilidad y consistencia de los 31 casos de prueba, la base de datos local (\texttt{version\_eco\_db}) debe encontrarse en un estado inicial controlado (\textit{Baseline}), el cual se establece mediante los scripts \texttt{01\_schema\_version\_eco.sql} y \texttt{02\_seed\_demo.sql}:

\begin{tcolorbox}[colback=lightgray,colframe=primary,title=\textbf{Estado Inicial de la Base de Datos (Prerrequisitos de Ejecución)}]
\small
\begin{enumerate}[leftmargin=0.6cm]
    \item \textbf{Sedes Registradas:}
    \begin{itemize}[noitemsep]
        \item Sede Central (administrada internamente por el sistema para producción).
        \item Sede A (ej: \textit{Sede Principal - Norte}) y Sede B (ej: \textit{Sede Chapinero}), ambas con inventario inicial $= 0$.
    \end{itemize}
    \item \textbf{Catálogo de Materias Primas de Referencia:}
    \begin{itemize}[noitemsep]
        \item Insumo A: \textit{Gel de Aloe Vera} (Unidad: \texttt{kg}, Nivel mínimo: 1.0).
        \item Insumo B: \textit{Aceite de Coco Orgánico} (Unidad: \texttt{L}, Nivel mínimo: 1.0).
        \item Insumo C: \textit{Glicerina Vegetal} (Unidad: \texttt{kg}, Nivel mínimo: 1.0).
        \item Insumo D (para pruebas de eliminación F1-CP-08): \textit{Sal Marina} ($5\,\text{kg}$, sin asociar a recetas).
    \end{itemize}
    \item \textbf{Catálogo de Productos y Fórmulas (BOM):}
    \begin{itemize}[noitemsep]
        \item Producto Patrón: \textit{Champú Vegano de Romero y Aloe (80g)}.
        \item Composición de la Fórmula (por 1 unidad): Insumo A = $0.5\,\text{kg}$, Insumo B = $0.2\,\text{L}$, Insumo C = $0.1\,\text{kg}$.
    \end{itemize}
    \item \textbf{Inventario Central Inicial:} 0 unidades de Champú Vegano en bodega central.
    \item \textbf{Ajuste Dinámico por Caso:} Antes de ejecutar cada corrida de F3, el stock disponible de A, B y C se ajusta a los valores exactos definidos en la columna \textit{Entradas} mediante el módulo de Materias Primas.
\end{enumerate}
\end{tcolorbox}

\subsection{Resumen de Casos Diseñados por Funcionalidad}
\begin{table}[htbp]
\centering
\small
\begin{tabular}{l c c c}
\toprule
\textbf{Funcionalidad} & \textbf{Total Casos} & \textbf{Técnicas Asignadas} & \textbf{Casos con Defecto Conocido} \\
\midrule
\textbf{F3: Registro de producción} & 12 & Arreglo Ortogonal $L_9(3^4)$ + Tabla de Decisión & 0 \\
\textbf{F4: Traslado de productos a sedes} & 7 & Tabla de Decisión + Valores Límite & 0 \\
\textbf{F1: Gestión de materia prima} & 8 & Tabla de Decisión + Valores Límite & 1 (F1-CP-03) \\
\textbf{F2: Creación de fórmulas} & 4 & Tabla de Decisión & 0 \\
\midrule
\textbf{TOTAL} & \textbf{31} & & \textbf{1} \\
\bottomrule
\end{tabular}
\caption{Consolidado del Alcance de Casos de Prueba Diseñados.}
\end{table}

\subsection{Precondiciones y Pasos de Ejecución Comunes (Procedimiento Operativo)}

\begin{tcolorbox}[colback=lightgray,colframe=secondary,title=\textbf{Procedimiento F3: Registro de Producción}]
\textbf{Precondiciones:}
\begin{enumerate}[noitemsep]
    \item Sistema desplegado y accesible desde el navegador web.
    \item Base de datos limpia con sede central, producto \textit{Champú vegano} y su fórmula ($A=0.5\,\text{kg}, B=0.2\,\text{L}, C=0.1\,\text{kg}$ por unidad), nivel mínimo 1 en cada materia prima.
    \item Inventario inicial del producto en sede central igual a 0.
    \item Stock inicial de A, B y C ajustado según las Entradas de cada caso.
\end{enumerate}
\textbf{Pasos de Ejecución:}
\begin{enumerate}[noitemsep]
    \item Ajustar el stock de A, B y C a los valores del caso desde el módulo \textit{Materias Primas}.
    \item Anotar los valores iniciales de A, B, C y del inventario del producto.
    \item Ingresar al módulo \textit{Producción}.
    \item Seleccionar el producto \textit{Champú vegano}.
    \item Ingresar la cantidad $Q$ indicada en el caso.
    \item Confirmar el registro de la producción.
    \item Consultar stock de A, B y C, inventario central y alertas, comparando con el resultado esperado.
\end{enumerate}
\end{tcolorbox}

\begin{tcolorbox}[colback=lightgray,colframe=secondary,title=\textbf{Procedimiento F4: Traslado de Productos a Sedes}]
\textbf{Precondiciones:}
\begin{enumerate}[noitemsep]
    \item Sistema desplegado y accesible desde el navegador.
    \item Producto \textit{Champú vegano} con inventario central = 20 (producido previamente).
    \item Sede A y Sede B registradas en el módulo \textit{Sedes}, ambas con inventario = 0.
\end{enumerate}
\textbf{Pasos de Ejecución:}
\begin{enumerate}[noitemsep]
    \item Anotar el inventario central y el de las sedes (módulos \textit{Distribución} e \textit{Inventario}).
    \item Ir al módulo \textit{Distribución} (se muestra formulario ``Enviar producto a sede'').
    \item Seleccionar el producto \textit{Champú vegano} y la sede destino.
    \item Ingresar la cantidad $Q$ y la referencia \texttt{"TEST"}.
    \item Pulsar el botón \textit{Distribuir}.
    \item Consultar el inventario central y el de las sedes para validar los nuevos saldos.
\end{enumerate}
\end{tcolorbox}

\begin{tcolorbox}[colback=lightgray,colframe=secondary,title=\textbf{Procedimiento F1 y F2: Materias Primas y Fórmulas}]
\textbf{F1 (Gestión de Materia Prima):}
\begin{enumerate}[noitemsep]
    \item Precondición: Existe la materia prima \textit{"Aloe vera"} ($10\,\text{kg}$) usada en fórmula, y \textit{"Sal"} ($5\,\text{kg}$) sin fórmula.
    \item Pasos: Ir a \textit{Materias Primas} $\rightarrow$ Ejecutar alta o eliminación con las entradas indicadas $\rightarrow$ Verificar listado, alertas e historial.
\end{enumerate}
\textbf{F2 (Creación de Fórmulas):}
\begin{enumerate}[noitemsep]
    \item Precondición: Existen insumos \textit{"Aloe vera"} y \textit{"Aceite de coco"} y el producto \textit{"Champú vegano"}.
    \item Pasos: Ir a \textit{Productos y Fórmulas} $\rightarrow$ Definir la fórmula con las materias y cantidades del caso $\rightarrow$ Consultar la fórmula guardada.
\end{enumerate}
\end{tcolorbox}

\begin{tcolorbox}[colback=primary!5,colframe=primary,title=\textbf{Nomenclatura y Variables de Prueba Utilizadas en la Matriz}]
Para facilitar la lectura y correspondencia con el diseño de Taguchi y Tablas de Decisión del equipo, se definen los siguientes parámetros operativos:
\begin{itemize}[noitemsep]
    \item \textbf{$Q$ (Quantity):} Cantidad numérica entera de unidades a producir (en F3) o a transferir (en F4).
    \item \textbf{$A$ (Materia Prima A):} \textit{"Aloe vera"} (Unidad: kg). Dosificación en fórmula: $0.5\,\text{kg}$ por unidad. Nivel mínimo de alerta: $1.0\,\text{kg}$.
    \item \textbf{$B$ (Materia Prima B):} \textit{"Aceite de coco"} (Unidad: kg). Dosificación en fórmula: $0.2\,\text{kg}$ por unidad. Nivel mínimo de alerta: $1.0\,\text{kg}$.
    \item \textbf{$C$ (Materia Prima C):} \textit{"Romero"} (Unidad: kg). Dosificación en fórmula: $0.1\,\text{kg}$ por unidad. Nivel mínimo de alerta: $1.0\,\text{kg}$.
    \item \textbf{Producto Evaluado:} \textit{"Champú vegano"}, cuya receta oficial vincula exactamente los insumos $A$, $B$ y $C$.
    \item \textbf{Criterio de Alerta Roja:} Se dispara automáticamente cuando el stock remanente de cualquier materia prima es menor o igual a su nivel mínimo ($1.0$).
\end{itemize}
\end{tcolorbox}

\subsection{Matriz Detallada de los 31 Casos de Prueba Diseñados}

\begin{longtable}{p{1.4cm} p{1.6cm} p{4.2cm} p{3.2cm} p{4.6cm}}
\caption{Matriz Completa de los 31 Casos de Prueba con Precondiciones y Entradas.}\label{tab:casos31}\\
\toprule
\textbf{ID Caso} & \textbf{Prior./Téc.} & \textbf{Precondición Específica} & \textbf{Entrada(s) de Prueba} & \textbf{Resultado Esperado} \\
\midrule
\endfirsthead
\multicolumn{5}{c}{\tablename\ \thetable\ -- \textit{Continuación de la página anterior}} \\
\toprule
\textbf{ID Caso} & \textbf{Prior./Téc.} & \textbf{Precondición Específica} & \textbf{Entrada(s) de Prueba} & \textbf{Resultado Esperado} \\
\midrule
\endhead
\midrule
\multicolumn{5}{r}{\textit{Continúa en la siguiente página\dots}} \\
\endfoot
\bottomrule
\endlastfoot
\multicolumn{5}{l}{\textbf{\color{primary}Módulo F3: Registro de Producción Centralizada}} \\
\midrule
\textbf{F3-CP-01} & Alta \newline AO 1 + TD & Fórmula de Champú activa ($A{=}0.5, B{=}0.2, C{=}0.1$). Stock: $A{=}10, B{=}7, C{=}6$. Prod. central: $0$. & $Q=10$ unid. a producir. & Producción exitosa. Descuento exacto: $A{=}5, B{=}5, C{=}5$. Stock champú $+10$. Sin alertas. \\
\textbf{F3-CP-02} & Alta \newline AO 2 + TD & Fórmula activa. Stock inicial: $A{=}17.5, B{=}5, C{=}2.5$. & $Q=25$ unid. & Producción registrada. Final: $A{=}5, B{=}0, C{=}0$. Inventario $+25$. Genera alertas rojas en B y C. \\
\textbf{F3-CP-03} & Crítica \newline AO 3 + TD & Fórmula activa. Stock inicial: $A{=}5.5, B{=}0.15, C{=}0.05$ (insuficiente en B y C). & $Q=1$ unid. & Rechazo por insuficiencia de B y C. Rollback total: stocks e inventario sin modificación. \\
\textbf{F3-CP-04} & Crítica \newline AO 4 + TD & Fórmula activa. Stock inicial: $A{=}5, B{=}2, C{=}0.95$ (requiere $1$ de C). & $Q=10$ unid. & Rechazo por falta de C. A y B siguen intactos en 5 y 2 (sin descuento parcial). \\
\textbf{F3-CP-05} & Crítica \newline AO 5 + TD & Fórmula activa. Stock inicial: $A{=}12.5, B{=}4.95, C{=}7.5$ (requiere $5$ de B). & $Q=25$ unid. & Rechazo por falta de B. A y C siguen en 12.5 y 7.5. Inventario sin alteraciones. \\
\textbf{F3-CP-06} & Alta \newline AO 6 + TD & Fórmula activa. Stock inicial: $A{=}0.5, B{=}5.2, C{=}0.1$. & $Q=1$ unid. & Producción exitosa. Final: $A{=}0, B{=}5, C{=}0$. Inventario $+1$. Alerta roja en A y C. \\
\textbf{F3-CP-07} & Crítica \newline AO 7 + TD & Fórmula activa. Stock inicial: $A{=}4.95, B{=}1.95, C{=}1$ (requiere $5$ de A y $2$ de B). & $Q=10$ unid. & Rechazo por falta de A y B. Stock de C permanece en 1. Sin creación de orden. \\
\textbf{F3-CP-08} & Crítica \newline AO 8 + TD & Fórmula activa. Stock inicial: $A{=}12.45, B{=}10, C{=}2.45$ (requiere $12.5$ de A y $2.5$ de C). & $Q=25$ unid. & Rechazo por falta de A y C. Insumo B sigue en 10. Inventario sin cambios. \\
\textbf{F3-CP-09} & Crítica \newline AO 9 + TD & Fórmula activa. Stock inicial: $A{=}0.45, B{=}0.2, C{=}5.1$ (requiere $0.5$ de A). & $Q=1$ unid. & Rechazo por falta de A. B y C siguen en 0.2 y 5.1. Transacción revertida. \\
\textbf{F3-CP-10} & Alta \newline TD (R1) & Producto con fórmula y materias primas disponibles en bodega. & $Q=0$ unid. & Rechazo por cantidad inválida ($Q \le 0$). Inventarios y stocks intactos. \\
\textbf{F3-CP-11} & Alta \newline TD (R1) & Producto con fórmula activa disponible. & $Q=-5$ unid. & Rechazo por valor negativo inválido. Bloqueo en validación frontal y backend. \\
\textbf{F3-CP-12} & Alta \newline TD (R1) & Materias primas con stock suficiente en bodega. & $Q=2.5$ unid. (decimal) & Rechazo por cantidad fraccionaria. Se requiere valor entero estricto. \\
\midrule
\multicolumn{5}{l}{\textbf{\color{primary}Módulo F4: Traslado de Productos a Sedes}} \\
\midrule
\textbf{F4-CP-01} & Alta \newline TD & Sede Central tiene 20 unid. Champú. Sede Norte tiene 0 unid. & Trasladar $Q=5$ unid. a Sede Norte. & Traslado exitoso. Central disminuye a 15, Sede Norte incrementa a 5. \\
\textbf{F4-CP-02} & Alta \newline TD + VL & Sede Central tiene 20 unid. Sede Norte tiene 0 unid. & Trasladar $Q=20$ unid. a Sede Norte (límite superior). & Traslado exitoso. Central queda en 0 unid.; Sede Norte queda en 20 unid. \\
\textbf{F4-CP-03} & Crítica \newline TD + VL & Sede Central tiene 20 unid. Sede Norte tiene 0 unid. & Trasladar $Q=21$ unid. a Sede Norte (límite $+ 1$). & Rechazo por stock insuficiente. Central mantiene 20 y Sede Norte mantiene 0. \\
\textbf{F4-CP-04} & Alta \newline TD & Sede Central cuenta con existencias disponibles ($>10$). & Trasladar $Q=2.5$ unid. a Sede Norte. & Rechazo por formato no entero. Bloqueo en control de entrada y validación API. \\
\textbf{F4-CP-05} & Alta \newline TD & Sede Central tiene 20 unid. Sede Sur tiene saldo previo de 3 unid. & Trasladar $Q=4$ unid. a Sede Sur. & Traslado exitoso acumulativo. Central queda en 13 unid., Sede Sur queda en 7 unid. \\
\textbf{F4-CP-06} & Alta \newline TD & Sedes origen y destino válidas y seleccionadas. & Trasladar $Q=0$ unid. a Sede Norte. & Rechazo por cantidad cero inválida ($Q \le 0$). Sin afectación en inventarios. \\
\textbf{F4-CP-07} & Alta \newline TD & Sedes origen y destino seleccionadas. & Trasladar $Q=-5$ unid. a Sede Norte. & Rechazo inmediato por valor negativo inválido. Saldos inalterados. \\
\midrule
\multicolumn{5}{l}{\textbf{\color{primary}Módulo F1: Gestión de Materia Prima}} \\
\midrule
\textbf{F1-CP-01} & Alta \newline TD & Insumo "Aceite de coco" no existe en el catálogo. & Nombre: "Aceite de coco", Unidad: "Litros", Stock: 20. & Registro exitoso. Insumo creado en catálogo con ID y stock inicial asignado. \\
\textbf{F1-CP-02} & Alta \newline TD & Existe previamente la materia prima "Aloe vera" en el sistema. & Nombre: "Aloe vera", Unidad: "kg", Stock: 5. & Rechazo por nombre duplicado exacto (Error 409). No se duplica en BD. \\
\textbf{F1-CP-03} & Alta \newline TD & Existe previamente la materia prima "Aloe vera" en el sistema. & Nombre: \texttt{"aloe vera "} (minúsculas y espacio final). & Rechazo por duplicado según RN-01 (normalización). \textit{[DEF-01: detectado fallo en backend].} \\
\textbf{F1-CP-04} & Alta \newline TD + VL & Catálogo disponible sin el insumo "Glicerina". & Nombre: "Glicerina", Unidad: "kg", Stock: 50 (límite superior). & Registro exitoso en el límite superior admitido de 50 kg/l (RN-04). \\
\textbf{F1-CP-05} & Alta \newline TD + VL & Catálogo disponible. Capacidad máxima definida en 50. & Nombre: "Glicerina", Unidad: "kg", Stock: 50.01 (límite $+ 0.01$). & Rechazo por exceder capacidad máxima de almacenamiento (RN-04). \\
\textbf{F1-CP-06} & Alta \newline TD & Formulario de materias primas abierto. & Nombre: "Glicerina", Unidad: "kg", Stock: -1. & Rechazo por cantidad negativa no permitida. Bloqueo de entrada. \\
\textbf{F1-CP-07} & Media \newline TD & Catálogo disponible. Insumo nuevo no registrado. & Nombre: "Glicerina", Unidad: "kg", Stock: 0. & Registro exitoso y generación inmediata de alerta visual roja por stock en 0 (RN-05). \\
\textbf{F1-CP-08} & Media \newline TD & "Aloe vera" está asociada a la fórmula activa de Champú. & Solicitar eliminación física de "Aloe vera". & Eliminación rechazada (Error 409) por restricción de integridad referencial (RN-02). \\
\midrule
\multicolumn{5}{l}{\textbf{\color{primary}Módulo F2: Creación de Fórmulas}} \\
\midrule
\textbf{F2-CP-01} & Alta \newline TD & Insumos Aloe y Aceite existen. Producto "Champú" sin receta. & Aloe vera (0.5 kg) + Aceite de coco (0.2 L). & Fórmula creada y vinculada exitosamente al producto. Consultable en catálogo. \\
\textbf{F2-CP-02} & Alta \newline TD & Producto registrado sin fórmula asociada. & Guardar fórmula sin agregar ningún insumo. & Rechazo: la receta exige al menos una materia prima (RN-07). \\
\textbf{F2-CP-03} & Alta \newline TD & Producto seleccionado; insumo "Aloe vera" ya agregado a la lista. & Agregar nuevamente "Aloe vera" (Aloe 0.5 + Aloe 0.3). & Rechazo por insumo duplicado dentro de la misma fórmula (RN-08). \\
\textbf{F2-CP-04} & Alta \newline TD & Producto y materia prima seleccionados en formulario. & Asignar cantidad requerida en cero ($Q{=}0$). & Rechazo: la cantidad requerida de insumo debe ser estrictamente $>0$. \\
\end{longtable}

% ==============================================================================
% 9. REPORTE DE RESULTADOS Y EVIDENCIAS DE EJECUCIÓN
% ==============================================================================
\section{Reporte de Resultados y Evidencias de Ejecución}

\subsection{Resumen Estadístico Consolidado}

Se ejecutó el ciclo completo de los \textbf{31 casos de prueba} sobre el software desplegado localmente en ambiente controlado:

\begin{table}[htbp]
\centering
\small
\begin{tabular}{l c c c c c r}
\toprule
\textbf{Feature / Módulo} & \textbf{Total} & \textbf{Aprobado} & \textbf{Fallido} & \textbf{Bloqueado} & \textbf{No Ejec.} & \textbf{\% Aprobado} \\
\midrule
\textbf{F3: Registro de producción} & 12 & 12 & 0 & 0 & 0 & \textbf{100.0\%} \\
\textbf{F4: Traslado de productos a sedes} & 7 & 7 & 0 & 0 & 0 & \textbf{100.0\%} \\
\textbf{F1: Gestión de materia prima} & 8 & 7 & 1 & 0 & 0 & \textbf{87.5\%} \\
\textbf{F2: Creación de fórmulas} & 4 & 4 & 0 & 0 & 0 & \textbf{100.0\%} \\
\midrule
\textbf{TOTAL GENERAL} & \textbf{31} & \textbf{30} & \textbf{1} & \textbf{0} & \textbf{0} & \textbf{96.77\%} \\
\bottomrule
\end{tabular}
\caption{Consolidado de Resultados de Ejecución de Pruebas.}
\end{table}

\begin{tcolorbox}[colback=accentgreen!10,colframe=accentgreen,title=\textbf{Dictamen de Certificación del Sistema}]
\textbf{Tasa Global de Aprobación:} \textbf{96.77\%} (30 casos aprobados de 31 ejecutados).\\
\textbf{Resultado:} \textbf{APROBADO PARA CERTIFICACIÓN}. Supera con holgura el umbral mínimo del 90\% establecido en el Plan Maestro de Pruebas, no presentándose ningún defecto de severidad crítica ni alta.
\end{tcolorbox}

\subsection{Matriz de Resultados de Ejecución Caso por Caso}

\begin{longtable}{p{1.6cm} p{3.4cm} p{7.8cm} p{2.2cm}}
\caption{Matriz Detallada de Ejecución de los 31 Casos de Prueba.}\label{tab:ejecucion31}\\
\toprule
\textbf{ID Caso} & \textbf{Resultado Esperado} & \textbf{Resultado Obtenido en Ejecución Real} & \textbf{Estado} \\
\midrule
\endfirsthead
\multicolumn{4}{c}{\tablename\ \thetable\ -- \textit{Continuación de la página anterior}} \\
\toprule
\textbf{ID Caso} & \textbf{Resultado Esperado} & \textbf{Resultado Obtenido en Ejecución Real} & \textbf{Estado} \\
\midrule
\endhead
\midrule
\multicolumn{4}{r}{\textit{Continúa en la siguiente página\dots}} \\
\endfoot
\bottomrule
\endlastfoot
\textbf{F3-CP-01} & Producción $+10$. Final: $A=5, B=5, C=5$. & Lote registrado. Saldos descontados exactamente. Inv $+10$. & \textbf{\textcolor{accentgreen}{Aprobado}} \\
\textbf{F3-CP-02} & Producción $+25$. Final: $A=5, B=0, C=0$. Alertas. & Lote creado. B y C en 0. Alertas rojas visibles en Dashboard. & \textbf{\textcolor{accentgreen}{Aprobado}} \\
\textbf{F3-CP-03} & Rechazo por falta de B y C. Stock intacto. & Excepción capturada. Stock de A, B y C sin cambios. & \textbf{\textcolor{accentgreen}{Aprobado}} \\
\textbf{F3-CP-04} & Rechazo por C. Sin descuento parcial de A y B. & Operación cancelada. A y B permanecen en 5 y 2. & \textbf{\textcolor{accentgreen}{Aprobado}} \\
\textbf{F3-CP-05} & Rechazo por B. A y C sin cambios. & Bloqueo exitoso. A y C conservan saldos $12.5$ y $7.5$. & \textbf{\textcolor{accentgreen}{Aprobado}} \\
\textbf{F3-CP-06} & Producción $+1$. Final: $A=0, B=5, C=0$. & Lote creado. Alertas de desabastecimiento en A y C. & \textbf{\textcolor{accentgreen}{Aprobado}} \\
\textbf{F3-CP-07} & Rechazo por falta de A y B. C intacto. & Bloqueo y Rollback confirmado. Sin alteraciones en BD. & \textbf{\textcolor{accentgreen}{Aprobado}} \\
\textbf{F3-CP-08} & Rechazo por falta de A y C. B intacto. & Bloqueo confirmado. B permanece en 10. & \textbf{\textcolor{accentgreen}{Aprobado}} \\
\textbf{F3-CP-09} & Rechazo por falta de A ($0.45 < 0.5$). & Bloqueo confirmado. Insumos B y C se mantienen intactos. & \textbf{\textcolor{accentgreen}{Aprobado}} \\
\textbf{F3-CP-10} & Rechazo por cantidad cero ($Q=0$). & Formulario bloquea el envío (control HTML5 \texttt{min="1"}). & \textbf{\textcolor{accentgreen}{Aprobado}} \\
\textbf{F3-CP-11} & Rechazo por cantidad negativa ($Q=-5$). & Validación de entrada rechaza el valor; sin cambios en BD. & \textbf{\textcolor{accentgreen}{Aprobado}} \\
\textbf{F3-CP-12} & Rechazo por cantidad decimal ($Q=2.5$). & El campo numérico entero impide registrar decimales. & \textbf{\textcolor{accentgreen}{Aprobado}} \\
\textbf{F4-CP-01} & Traslado 5: Central 15, Sede A 5. & Despacho exitoso. Central queda en 15 y Sede A en 5. & \textbf{\textcolor{accentgreen}{Aprobado}} \\
\textbf{F4-CP-02} & Traslado límite 20: Central 0, Sede A 20. & Despacho exitoso. Central queda exactamente en 0. & \textbf{\textcolor{accentgreen}{Aprobado}} \\
\textbf{F4-CP-03} & Traslado 21 excede stock: Rechazo. & Mensaje en pantalla: ``Stock insuficiente en central''. & \textbf{\textcolor{accentgreen}{Aprobado}} \\
\textbf{F4-CP-04} & Traslado decimal $2.5$: Rechazo. & Formulario bloquea el valor fraccionario. & \textbf{\textcolor{accentgreen}{Aprobado}} \\
\textbf{F4-CP-05} & Suma acumulada en Sede B: Central 13, Sede 7. & Saldos finales coinciden exactamente con la suma teórica. & \textbf{\textcolor{accentgreen}{Aprobado}} \\
\textbf{F4-CP-06} & Traslado cantidad 0: Rechazo. & Validación de formulario impide enviar cantidad nula. & \textbf{\textcolor{accentgreen}{Aprobado}} \\
\textbf{F4-CP-07} & Traslado negativo $-5$: Rechazo. & Bloqueo en interfaz y rechazo en API. & \textbf{\textcolor{accentgreen}{Aprobado}} \\
\textbf{F1-CP-01} & Alta válida: Aceite de coco registrado. & Insumo insertado exitosamente en catálogo. & \textbf{\textcolor{accentgreen}{Aprobado}} \\
\textbf{F1-CP-02} & Nombre duplicado exacto: Rechazo. & Retorna error HTTP 409 por violación de unicidad. & \textbf{\textcolor{accentgreen}{Aprobado}} \\
\textbf{F1-CP-03} & Duplicado con mayúsculas y espacios: Rechazo. & \textbf{El sistema permitió crear ``aloe vera '' habiendo ``Aloe vera''.} & \textbf{\textcolor{accentred}{Fallido}} \\
\textbf{F1-CP-04} & Cantidad en tope 50: Registro válido. & Registrado con $50\,\text{kg}$ en almacén. & \textbf{\textcolor{accentgreen}{Aprobado}} \\
\textbf{F1-CP-05} & Cantidad $50.01$: Rechazo por exceder 50. & Operación bloqueada; no permite superar el límite. & \textbf{\textcolor{accentgreen}{Aprobado}} \\
\textbf{F1-CP-06} & Cantidad negativa $-1$: Rechazo. & Control de formulario impide valores negativos. & \textbf{\textcolor{accentgreen}{Aprobado}} \\
\textbf{F1-CP-07} & Cantidad 0: Registrado con alerta. & Insumo registrado con 0; semáforo rojo se activa en Dashboard. & \textbf{\textcolor{accentgreen}{Aprobado}} \\
\textbf{F1-CP-08} & Eliminar insumo en fórmula: Bloqueado. & Mensaje de error 409: insumo con receta activa. & \textbf{\textcolor{accentgreen}{Aprobado}} \\
\textbf{F2-CP-01} & Fórmula válida: Receta guardada. & Composición almacenada y consultable en producto. & \textbf{\textcolor{accentgreen}{Aprobado}} \\
\textbf{F2-CP-02} & Fórmula sin insumos: Rechazo. & Sistema exige vincular al menos una materia prima. & \textbf{\textcolor{accentgreen}{Aprobado}} \\
\textbf{F2-CP-03} & Materia repetida: Rechazo. & Interfaz impide agregar dos veces el mismo ingrediente. & \textbf{\textcolor{accentgreen}{Aprobado}} \\
\textbf{F2-CP-04} & Cantidad en 0: Rechazo. & Se exige cantidad requerida $>0$. & \textbf{\textcolor{accentgreen}{Aprobado}} \\
\end{longtable}

\subsection{Documentación del Defecto Encontrado}

\begin{tcolorbox}[colback=lightgray,colframe=accentred,title=\textbf{Ficha de Defecto: DEF-01 (F1-CP-03 -- Violación de RN-01)}]
\small
\textbf{Identificador:} DEF-01 \quad \textbf{Severidad:} Media \quad \textbf{Prioridad:} Media \quad \textbf{Estado:} Abierto\\
\textbf{Funcionalidad Afectada:} F1 -- Gestión de Materia Prima / Endpoint \texttt{POST /api/materias-primas}\\[2mm]
\textbf{Descripción:}\\
Al crear una nueva materia prima cuyo nombre coincide con una previamente registrada variando únicamente en mayúsculas o espacios terminales (ejemplo: ingresar \texttt{"aloe vera "} cuando ya existe \texttt{"Aloe vera"}), el sistema acepta la inserción con código HTTP 201 en lugar de rechazarla con HTTP 409 Conflict.\\[2mm]
\textbf{Pasos para Reproducir:}
\begin{enumerate}[noitemsep]
    \item Acceder al módulo \textit{Materias Primas}.
    \item Diligenciar el formulario con Nombre: \texttt{"aloe vera "}, Unidad: \texttt{kg}, Cantidad: \texttt{5}.
    \item Presionar el botón \textit{Guardar}.
    \item Observar que se crea un segundo registro redundante en el listado.
\end{enumerate}
\textbf{Causa Raíz:} El backend no aplica la función de sanitización \texttt{.trim().toLowerCase()} antes de ejecutar la consulta de verificación ni en la cláusula de inserción SQL.\\
\textbf{Solución Recomendada:} Normalizar la entrada con \texttt{nombre.trim().toLowerCase()} en el controlador y asegurar el índice único relacional:\\
\texttt{CREATE UNIQUE INDEX uq\_mp\_norm ON materia\_prima (LOWER(TRIM(nombre))) WHERE deleted\_at IS NULL;}
\end{tcolorbox}

% ==============================================================================
% 10. CONCLUSIONES
% ==============================================================================
\section{Conclusiones y Evaluación de la Calidad}

\begin{enumerate}
    \item \textbf{Efectividad de la Estrategia de Pruebas:} La combinación de Arreglos Ortogonales ($L_9$) y Tablas de Decisión permitió someter a prueba exhaustiva la funcionalidad crítica de producción con solo 12 corridas, demostrando que el procedimiento almacenado de PostgreSQL cumple rigurosamente la regla de atomicidad (RN-10).
    \item \textbf{Aprobación Formal del Software:} Al alcanzar un \textbf{96.77\% de casos superados} frente al estándar del 90\%, el sistema Version-Eco se certifica como estable, consistente y apto para su utilización operativa.
    \item \textbf{Identificación Exitosa de Defectos:} Las pruebas manuales planificadas permitieron capturar y aislar el defecto conocido DEF-01 en el catálogo de materias primas, proporcionando la solución técnica preventiva para el siguiente sprint.
\end{enumerate}

\end{document}
